% =========================================================
% Stern–Brocot Limit Theory — Supporting Results
% For insertion into: volume-ii/rationals/notes/
% These results support and justify Theorem 10.102
% (Stern–Brocot Limit Theorem).
% =========================================================

% ---------------------------------------------------------
% Toolkit
% ---------------------------------------------------------

\begin{tcolorbox}[colback=gray!6, colframe=gray!40, arc=2pt,
  left=6pt, right=6pt, top=4pt, bottom=4pt,
  title={\small\textbf{Stern–Brocot Limit Theory — Quick Reference}},
  fonttitle=\small\bfseries]
\small
\begin{tabular}{l l l l}
\toprule
\textbf{Label} & \textbf{Statement} & \textbf{Proof method} & \textbf{Detail} \\
\midrule
L10.A & Denominators grow on any infinite path        & Induction + det.\ identity      & \hyperref[lem:sb-denom-growth]{$\downarrow$} \\
L10.B & Interval lengths shrink to zero               & Mediant Gap Lemma + L10.A       & \hyperref[lem:sb-length-zero]{$\downarrow$} \\
P10.A & Nested SB intervals contain a unique point    & Completeness of $\mathbb{R}$    & \hyperref[prop:sb-unique-point]{$\downarrow$} \\
C10.A & Every infinite path determines a unique limit & L10.B + P10.A                   & \hyperref[cor:sb-path-limit]{$\downarrow$} \\
C10.B & The limit of an infinite path is irrational   & Contrapositive: rational $\Rightarrow$ finite path & \hyperref[cor:sb-limit-irrational]{$\downarrow$} \\
\bottomrule
\end{tabular}
\end{tcolorbox}

% ---------------------------------------------------------
% Lemma A — Denominators grow
% ---------------------------------------------------------

\begin{lemma}[Denominator growth on infinite Stern--Brocot paths]
\label{lem:sb-denom-growth}
Let $a_0/b_0, a_1/b_1, a_2/b_2, \ldots$ be the sequence of fractions
encountered along any infinite path in the Stern--Brocot tree (i.e.\ the
sequence of interval-endpoint fractions at successive levels). Then
\[
  b_0 < b_1 < b_2 < \cdots,
\]
and in particular $b_n \to \infty$.
\end{lemma}

\begin{remark}[Fully quantified form]
$\forall n \in \mathbb{N},\ b_n < b_{n+1}$.
In particular, $\forall M > 0\ \exists N \in \mathbb{N}\ \forall n \ge N,\ b_n > M$.
\end{remark}

\begin{remark}[Intuition]
At each step in the tree, the new fraction produced by the mediant has
denominator equal to the \emph{sum} of the two parent denominators.
Since both parents have positive denominators, the new denominator
strictly exceeds each of them. A path that turns left or right at each
level therefore visits fractions with strictly increasing denominators.
\end{remark}

\begin{remark}[Proof strategy]
By induction on the depth of the path. At each step one endpoint
is replaced by the mediant; its denominator is the sum of the two
current denominators, which strictly exceeds the replaced endpoint's
denominator. The unboundedness follows because a strictly increasing
sequence of positive integers is unbounded.
\end{remark}

% ---------------------------------------------------------
% Lemma B — Interval lengths shrink to zero
% ---------------------------------------------------------

\begin{lemma}[Stern--Brocot interval lengths shrink to zero]
\label{lem:sb-length-zero}
Let $(I_n)_{n \ge 0}$ be the sequence of nested rational intervals generated
by the Stern--Brocot refinement process targeting a real number $x > 0$,
where $I_n = [a_n/b_n,\; c_n/d_n]$.  Write $\ell_n = c_n/d_n - a_n/b_n$
for the length of $I_n$.  Then $\ell_n \to 0$.
\end{lemma}

\begin{remark}[Fully quantified form]
$\forall \varepsilon > 0\ \exists N \in \mathbb{N}\ \forall n \ge N,\ \ell_n < \varepsilon$.
\end{remark}

\begin{remark}[Intuition]
By the Mediant Gap Lemma, $\ell_n = 1/(b_n d_n)$.
By Lemma~\ref{lem:sb-denom-growth}, both $b_n$ and $d_n$ tend to infinity,
so $1/(b_n d_n) \to 0$.
\end{remark}

\begin{remark}[Proof strategy]
The Farey gap identity gives $\ell_n = 1/(b_n d_n)$ (this is the content
of the Gap Between Farey Neighbors, Theorem~10.65, applied to each
successive Stern--Brocot interval). Lemma~\ref{lem:sb-denom-growth} gives
$b_n, d_n \to \infty$; therefore $1/(b_n d_n) \to 0$.
The formal step uses the Archimedean property: given $\varepsilon > 0$,
choose $N$ so that $b_N d_N > 1/\varepsilon$.
\end{remark}

\begin{remark}[Dependency note]
This lemma requires the Farey gap identity in the form $c/d - a/b = 1/(bd)$
for Farey neighbors.  In the Stern--Brocot construction the two endpoints
of $I_n$ are always Farey neighbors (this is the determinant condition
$b_n d_n' - a_n c_n' = 1$ propagated by mediant insertion), so the
identity applies at every level.
\end{remark}

% ---------------------------------------------------------
% Proposition A — Nested intervals contain a unique point
% ---------------------------------------------------------

\begin{proposition}[Nested Stern--Brocot intervals contain a unique limit point]
\label{prop:sb-unique-point}
Let $(I_n)_{n \ge 0}$ be the nested sequence of closed rational intervals
generated by the Stern--Brocot refinement process targeting a real number
$x > 0$, with $I_n = [a_n/b_n,\; c_n/d_n]$.  Then there exists a unique
real number $\xi \in \bigcap_{n=0}^{\infty} I_n$.
\end{proposition}

\begin{remark}[Fully quantified form]
$\exists!\, \xi \in \mathbb{R}$ such that
$\forall n \in \mathbb{N},\ a_n/b_n \le \xi \le c_n/d_n$.
\end{remark}

\begin{remark}[Intuition]
The sequence of left endpoints $(a_n/b_n)$ is increasing and bounded above
by $c_0/d_0$; the sequence of right endpoints $(c_n/d_n)$ is decreasing and
bounded below by $a_0/b_0$.  Both sequences converge in $\mathbb{R}$
(by the monotone convergence theorem), and their limits must coincide
because the interval lengths tend to zero.
\end{remark}

\begin{remark}[Proof strategy]
Show that the left endpoint sequence is monotone increasing and bounded
above; the right endpoint sequence is monotone decreasing and bounded
below.  Apply the Monotone Convergence Theorem in $\mathbb{R}$ to both.
Then use $\ell_n \to 0$ (Lemma~\ref{lem:sb-length-zero}) to conclude
the two limits are equal.  Uniqueness follows: any two points in
$\bigcap I_n$ would differ by at most $\ell_n$ for every $n$, hence
by zero.
\end{remark}

\begin{remark}[Where completeness enters]
The Monotone Convergence Theorem in $\mathbb{R}$ is a consequence of the
least-upper-bound property, which $\mathbb{Q}$ does not satisfy.  This is
precisely why the limit $\xi$ is a real number and need not be rational.
The result would fail if we worked only in $\mathbb{Q}$.
\end{remark}

% ---------------------------------------------------------
% Corollary A — Every infinite path has a unique limit
% ---------------------------------------------------------

\begin{corollary}[Every infinite Stern--Brocot path converges]
\label{cor:sb-path-limit}
Let $x > 0$ be a real number and let $(I_n)$ be the Stern--Brocot
refinement sequence targeting $x$.  Then the sequence of left endpoints
$(a_n/b_n)$ and the sequence of right endpoints $(c_n/d_n)$ both
converge in $\mathbb{R}$ to a common limit $\xi$, and $\xi = x$.
\end{corollary}

\begin{remark}[Fully quantified form]
$\displaystyle\lim_{n \to \infty} \frac{a_n}{b_n}
 = \lim_{n \to \infty} \frac{c_n}{d_n} = x$.
\end{remark}

\begin{remark}[Why $\xi = x$]
By construction, $x \in I_n$ for every $n$ (we always retain $x$ in the
interval when choosing which endpoint to replace by the mediant).
Therefore $x \in \bigcap I_n$. By Proposition~\ref{prop:sb-unique-point}
the intersection contains exactly one point, so $\xi = x$.
\end{remark}

\begin{remark}[Proof strategy]
This is an immediate consequence of Proposition~\ref{prop:sb-unique-point}
together with the observation that $x \in I_n$ for all $n$ by the
construction rule (replace whichever endpoint lies on the same side as
the mediant, keeping $x$ in the interval at every step).
\end{remark}

% ---------------------------------------------------------
% Corollary B — The limit of an infinite path is irrational
% ---------------------------------------------------------

\begin{corollary}[Infinite Stern--Brocot paths converge to irrationals]
\label{cor:sb-limit-irrational}
A positive real number $x$ corresponds to a finite path in the
Stern--Brocot tree if and only if $x$ is rational.  Equivalently,
$x$ corresponds to an infinite path if and only if $x$ is irrational.
\end{corollary}

\begin{remark}[Fully quantified form]
$x \in \mathbb{Q}_{>0}\ \Longleftrightarrow\ $ the Stern--Brocot refinement
process for $x$ terminates in finitely many steps.
Equivalently:
$x \notin \mathbb{Q}\ \Longleftrightarrow\ $ the refinement process
runs for infinitely many steps.
\end{remark}

\begin{remark}[Intuition]
The mediant is inserted and tested against $x$ at each step.  The process
terminates exactly when $x$ equals the mediant, which forces $x$ to be
rational (and in fact equal to the mediant, which is a rational number in
reduced form).  If $x$ is irrational it can never equal a mediant, so the
process never terminates.
\end{remark}

\begin{remark}[Proof strategy]
Prove both directions separately.
$(\Rightarrow)$: If $x = p/q$ is rational in reduced form, show that the
path reaches $x$ in finitely many steps via a descent argument analogous to
the Euclidean algorithm (the denominators of successive mediants decrease by
a bounded amount and must eventually reach $q$).
$(\Leftarrow)$: If the path is finite, it terminates at a mediant, which is
rational. Contapositive: if $x$ is irrational, the path cannot terminate.
\end{remark}